\chapter {Considerações Finais}

Neste capítulo é apresentado um resumo dos objetivos, desafios e soluções deste trabalho.
Na Seção \ref{sec:atividades-e-resultados} é apresentado o resumo da obra, e na Seção \ref{sec:trabalhos-futuros} são discutidos potenciais pontos de atuação para trabalhos futuros.

\section {Atividades e Resultados}
\label{sec:atividades-e-resultados}

Este trabalho teve como objetivo demonstrar uma jornada de desenvolvimento de funcionalidades levando em consideração requisitos levantados por \textit{stakeholders} e outros times, tendo como foco principal o desenvolvimento de funcionalidades com requisitos de desempenho mais rigorosos.

Com o surgimento de problemas na plataforma Brasil Participativo em 2023 e com a análise realizada pela Dataprev, um plano de atividades foi traçado considerando também as necessidades levantadas pelos \textit{stakeholders}.
As funcionalidades de textos participativos e cadastro de usuários foram tomadas como objetos de estudo e de melhorias.

A metodologia adotada se baseou no desenvolvimento contínuo.
Uma série de requisitos iniciais foram definidos, e em cada ciclo de desenvolvimento novos requisitos eram priorizados para serem trabalhados. Em suma, os requisitos foram:

\begin{itemize}
	\item Diminuir o tempo de renderização da página de visualização e de edição do texto participativo;
	\item Possibilitar a criação de textos participativos através de um editor de textos próprio do Brasil Participativo, contando com o uso do comando copiar e colar, trazendo texto de fontes como o Microsoft Word e Google Docs;
	\item Refatoração da página de visualização do rascunho (edição de parágrafos), trazendo-a para o mesmo padrão de \textit{design} que a tela de visualização;
	\item Refatorar a jornada de criação de textos participativos de acordo com o acordado com o time de produto;
  \item Resolver o problema de lentidão e indisponibilidade da funcionalidade de cadastro de usuários exibido durante o PPA.
\end{itemize}

Durante a jornada de desenvolvimento, uma série de modificações foram realizadas desde as páginas de exibição de conteúdo do texto até mesmo na forma como o usuário administrador submete o material na ferramenta.

Na primeira \textit{sprint} foi posta à prova as principais ideias de solução, removendo o uso indiscriminado de \textit{partials} do \textit{frontend}, e adicionando um editor de textos \textit{rich text} para envio de texto.
Juntamente, foi criado um algoritmo de \textit{parser}, a fim de separar parágrafos, de tal forma que cada um deles seja exibido e receba interações de forma isolada.
Melhorias de interface também foram conduzidas para que a ferramenta seguisse o padrão de estilo \texttt{Gov.Br} já adotado na plataforma.

Na segunda e terceira \textit{sprint} houve ajustes de usabilidade, aparência e também a correção de defeitos no \textit{parser} da ferramenta.
Pode-se dizer que esses dois ciclos foram utilizados para polir o que foi entregue já no primeiro.

Na quarta \textit{sprint}, surgiu a necessidade de implementar de forma diferente a edição e criação de textos, de tal maneira que fosse possível manipular tabelas nativamente na plataforma.
O editor de textos inicialmente disponibilizado, o Trix, foi substituído pelo editor de textos Jodit; e com isso foi também necessária a substituição do algoritmo de \textit{parser}.

E, por fim, na quinta \textit{sprint} foi abordado o problema de lentidão na criação de novos usuários, realizando ajustes em um ponto específico: a desambiguação de \textit{nicknames}, melhorando em 25\% o desempenho da consulta raiz do problema.

Ao final, foram conduzidos testes de carga para avaliar os ganhos de desempenho com as alterações propostas no trabalho.
Os resultados foram promissores, com reduções de tempo gasto de até 96\%.

\section {Trabalhos Futuros}
\label{sec:trabalhos-futuros}

No ponto em que se encerraram as atividades mencionadas neste trabalho, nota-se que existem diversas aberturas para melhorias, correções e evoluções da plataforma e das funcionalidades trabalhadas.
De forma não exaustiva as oportunidades de atuação são discutidas a seguir.

\subsection {Texto Participativo}

Com o uso do Jodit para captura do \textit{input} do usuário, existe a possibilidade, e também necessidade de adicionar validação sobre o conteúdo a ser processado e guardado no banco de dados.
O texto submetido pelo usuário é enviado no formato HTML, e dado que essa linguagem permite, além de estruturar textos, especificar \textit{scripts} e estilo em seu conteúdo, é necessário que seja feito um pós-processamento no lado \textit{backend}.
O Jodit foi configurado de forma que atributos de estilo e \textit{scripts} sejam removidos após a colagem de textos advindos de fontes externas.
Mas, nada impede que um usuário mal-intencionado sequestre uma conta de um usuário administrador para submeter HTML com conteúdo malicioso.

Com a remoção de \textit{tags} de estilo e \textit{script} do texto, algumas inconsistências surgem ao copiar e colar textos vindos do Google Docs, por exemplo.
Após observar o comportamento do copiar durante os testes, notou-se que a estilização do texto, como negrito, itálico e espaçamento é feita com o uso das \textit{tags} de estilo.
Isto é, um texto negrito, por exemplo, ao invés de ser representado através de uma \textit{tag} HTML, era representado através de atributos CSS.
Para esses casos, o conteúdo colado trazia o texto inteiro envolto sob uma \textit{tag} \texttt{<strong>}, dando a impressão de que tudo está em negrito, mesmo que no Google Docs isso não fosse o caso.

Na perspectiva de segurança, existe também a necessidade da criação de um mecanismo inteligente para validação das imagens submetidas.
Hoje as imagens são enviadas \textit{inline} no corpo do texto no formato \texttt{Data URL}.
Uma vez que o texto é processado pelo \textit{backend}, essas imagens são transformadas em \texttt{blobs} e enviadas ao \texttt{ActiveStorage} sem muita preocupação com seu conteúdo e tamanho.

Ainda sobre imagens, uma possível melhoria seria a criação de um mecanismo de envio assíncrono para o \textit{backend}.
Com uma \textit{controller} dedicada para o processamento dessas imagens, as rotinas de validação e persistência aconteceriam de forma granular e independente, removendo do \textit{parser} um pouco da sua complexidade.
Para que essa melhoria seja efetiva, também é necessário que seja implementado um mecanismo de limpeza do \texttt{storage}, de tal forma que imagens órfãs sejam removidas posteriormente, dado que uma imagem pode ser submetida e apagada logo em seguida sem nem mesmo persistir o texto no banco de dados.

A criação de textos grandes ainda sofre com o problema de lentidão, uma vez que é necessário separar cada parágrafo em um objeto único para ser salvo no formato de proposta no banco de dados.
Esses objetos possuem relacionamentos com outras entidades e também levam à criação de outros objetos para guardar metadados de autoria e auditoria.
Visando o ganho de desempenho, idealmente cada parágrafo do texto não deveria utilizar o modelo de propostas para operar, mas sim um modelo próprio, ou uma estrutura de dados que permita sua manipulação e exibição de forma rápida.

\subsection {Cadastro de Usuários}

No cadastro de usuários, com a troca do \texttt{ILIKE} pelo uso da função \texttt{LOWER}, observou-se uma melhora de desempenho.
Porém, existem formas atingir resultados ainda melhores na desambiguação de \textit{nicknames} através do uso de índices ou removendo a função \texttt{LOWER} da consulta.
A adoção de um índice poderia melhorar o desempenho da consulta, garantindo que a função \texttt{LOWER} não precise ser executada para cada tupla.
Outra forma seria remover o uso do \texttt{LOWER}, mas traria consigo uma série de cuidados necessários: garantir que os \textit{nicknames} são submetidos ao banco sempre em minúsculo, e os \textit{nicknames} já existentes estejam em minúsculo.

\subsection {Vazamento de Memória e Indisponibilidade do Banco de Dados}

Pela análise da Dataprev, foi visto que o vazamento de memória é um problema real, e que exige esforço do time de infraestrutura para corrigir essa falha de tempos em tempos.
Outro problema constatado foi a indisponibilidade do banco em diversos momentos de alta carga de acesso de usuários.
Existem ferramentas que podem ser adotadas na camada do \textit{backend} para potencialmente mitigar esses problemas.

Para melhorar o sistema de \textit{pool} do banco, pode-se explorar a utilização de ferramentas como o \texttt{pg-bouncer}.
E, para aprimorar o uso de memória, pode-se examinar o uso de outros mecanismos de alocação de memória, como, por exemplo, o \texttt{jemalloc}.